\section{The phase diagram}
\label{sec:results}

We simulate societies of $N=40$ agents over a $200\times200$ grid in $(d,\fd)$, one
society per grid point, and measure the five order parameters of
\S\ref{sec:order-params} after a fixed number of interactions, with both classes
discriminating (case~\eqref{eq:case6}). Parameters and their provenance are in
Appendix~\ref{app:parameters}; each pixel is a single realisation, so the speckle
is society-to-society fluctuation, not a rendering artefact.

\begin{figure}[t]
  \centering
  \includegraphics[width=\linewidth]{figures/order_parameter_maps}
  \caption{\textbf{All five order parameters over the $(d,\fd)$ plane}, for a
  simple agenda ($\alpha=0.17$, top) and a complex one ($\alpha=3.3$, bottom).
  $\Rmuc$ is the discrimination order parameter, flipping sign at $d\approx0$ ever
  more sharply as $\fd$ grows; $\Rcw$ is non-zero only for $d>0$ and, for the
  simple agenda, \emph{falls back} at the largest $d$, which is what separates the
  two discriminatory phases; $\BI$ separates the agenda sizes, featureless above
  and saturated below; and $\BA$ is the only panel in which the frustrated region
  at lower left announces itself, going strongly negative where everything else
  merely goes to zero.}
  \label{fig:correlations}
\end{figure}

\paragraph{Discrimination needs a quorum, and then arrives abruptly.} At small
$\fd$ the society stays neutral no matter how strong the bias: the few
discriminating agents are outvoted by the class-blind majority they learn from.
The trust--class correlation at large $d$ first exceeds $0.5$ around
$\fd\approx0.4$ ($0.43$ for the simple agenda, $0.35$ for the complex one), and past
that the transition in $d$ sharpens quickly: at $\fd=1$ the band in which
$|\Rmuc|<0.25$ is only $0.54$ wide, and $\Rmuc$ runs from $-0.99$ at $d=-1$ to
$+0.99$ at $d=+1$. This is the first practical consequence of the model: the state
of the population is not a smooth function of how biased its members are, and a
minority of biased agents leaves no trace at all until it reaches a threshold
fraction.

\paragraph{Two discriminatory phases, and only the simple agenda has both.} Both
discriminatory regions have distrust aligned to class; they differ in whether
opinion is aligned too. Along $\fd=0.9$ in the top row of
Figure~\ref{fig:correlations}, $\Rcw$ rises to $0.44$ at $d\approx0.6$ and then
\emph{falls} to $0.30$ by $d=1$, while $\Rmuc$ keeps climbing, from $0.78$ to
$0.87$ --- a gap of $0.55$ where the two started level. Ideological balance follows
$\Rcw$ down, $\BI$ dropping from $0.19$ to $0.06$ across the same row. At moderate bias the society splits into two camps
that both distrust and disagree with each other (region~III); at strong bias the
disagreement becomes unnecessary and class membership alone carries the hostility
(region~IV), which is what the magenta corner of Figure~\ref{fig:phase} shows.

The complex agenda has no region~IV. In the bottom row, $\Rmuc$ and $\Rcw$ track
each other across the whole $d>0$ half --- $0.84$ against $0.85$ at $d=0.15$, $0.91$
against $0.91$ at $d=0.75$, and still only $0.04$ apart at $d=1$ against the simple
agenda's $0.55$ --- while $\BI$ saturates at $0.98$. When there are many issues to disagree about, the ideological split
tracks the class split completely, so distrust and disagreement never come apart.
Class-only hostility is available only to a society whose agenda is too narrow for
ideology to reorganise around the class boundary.

\paragraph{Favouring the out-group produces frustration, not inverted
discrimination.} For $d<0$ the correlations are the sign-mirror of the $d>0$ side,
which invites reading the left half of the diagram as ``reverse discrimination''.
The balance measures say otherwise (Figure~\ref{fig:correlations}). In the
discriminatory region $\BA=+0.97$: nearly every triple is affectively balanced, as
two coherent mutually hostile blocs require. At $d=-1$, $\fd=0.9$,
$\BA=-0.71$ for both agendas --- strongly \emph{anti}-balanced --- with
$\BI=0.00$. No consistent assignment of trust exists:
each agent is pulled towards the other class while its own neighbours are pulled
towards theirs, and the society settles into a spin glass rather than into a
mirror-image society. This is also the one place where $\Rwmu$ collapses, to
$0.03$. Pair correlations alone would report the region as simply inverted; the
triple-level measurement identifies it as disordered.

\begin{figure}[t]
  \centering
  \begin{subfigure}{0.43\linewidth}
    \includegraphics[width=\linewidth]{figures/phase_diagram}
    \caption{The four phases.}
    \label{fig:phase}
  \end{subfigure}\hfill
  \begin{subfigure}{0.49\linewidth}
    \includegraphics[width=\linewidth]{figures/agenda_trajectories}
    \caption{Which polarisation comes first.}
    \label{fig:trajectories}
  \end{subfigure}
  \caption{\textbf{(a) Four collective states.} The three correlations of
  Figure~\ref{fig:correlations} (simple agenda) composited as colour channels ---
  red $\Rmuc$, green $\Rcw$, blue $\Rwmu$ --- so colour names the phase: (I) dark,
  frustrated reverse discrimination; (II) blue-violet, neutral polarisation
  unrelated to class; (III) pale, ideology tracking class; (IV) magenta, class
  alone. \textbf{(b) Balance trajectories} at $d=0$, one curve per agenda
  complexity, markers equally spaced in time (they crowd at the end because the
  annealed dynamics slows). Simple agendas run above the diagonal, complex ones
  below, crossing over at $\alpha\approx1.7$.}
  \label{fig:phase-and-traj}
\end{figure}

\paragraph{Which polarisation comes first depends on the agenda.} With no
discrimination at all, a society still polarises. Figure~\ref{fig:trajectories}
follows its path through the $(\BI,\BA)$ plane from the random initial condition,
where half the triples are frustrated, towards the frustration-free corner. For a
simple agenda the path runs up the $\BA$ axis: affective structure forms almost
completely before ideological structure begins, ending at
$(\BI,\BA)=(0.02,0.97)$ for $\alpha=0.03$. For a complex agenda the path stays
below the diagonal, ending at $(0.97,0.74)$ for $\alpha=333$: opinions separate
first and distrust follows. The crossover is at $\alpha\approx1.7$, where the
trajectory lies along the diagonal --- that is, when the number of issues in play
is comparable to the dimension of the space they are represented in.

Read as a mechanism, this says that the question of whether groups come to dislike
each other because they disagree, or come to disagree because they dislike each
other, has a parameter attached to it rather than an answer. Where the agenda is
narrow --- few issues, functioning as symbols --- affect leads. Where it is broad,
ideology leads. The same parameter explains why region~IV exists only in the top
row of Figure~\ref{fig:correlations}: with a simple agenda and a strong field,
ideology has not had time to reorganise around the class boundary, so hostility
arrives without it.
